% Pursuing Stacks (cotes 134-2 to 134-8): archive cote 134-7, site page 2 (PDF page 3), typescript page 444.
% typescript. Draft compared: ps6.tex:46-100.
% AI-transcribed by Claude Opus 5.5 (claude-opus-5-5, adaptive thinking, effort high), Message Batches API, batch ps3 (msgbatch_01XdFmuVVBAcgc7EFxdtCqsj), 2026-11-12; not yet checked by eye.
\scanpage{134-7/2}
\mspage{444}

truncation involving only the higher homotopy groups ($\hand{\pi}_i$ with $i\geqslant 2$).
This suggests that for any inte\xout{\unsure{x}}ger $n\geqslant 1$, there may be a ``schematization
% NOTE: the right edge of the scan cuts this line after "schematization"; the closing quotation marks are not visible
above level n'' for a given homotopy type, leaving the Cartan-Serre
truncation of \xout{\unsure{dimension}} level n discrete, and ``schematizing'' the
corresponding \inserted{total} Postnikov fiber (involving homotopy groups $\hand{\pi}_i$ for $i\geqslant n$).
Maybe such a schematization can be constructed equally for level 0 ,
even without assuming the fundamental groupoid to be abelian, \xout{by replac}
only nilpotent - but then we may have to change from ground ring $\underline{Z}$ to
the considerably coarser one $\underline{Q}$ (compare comments at the end of section
94). In any case, the key step in this approach would consist in \xout{the}
checking that, after schematization has been carried through succesful-%
ly up to a certain level in the successive elementary Postnikov fibrati-%
ons, the next elementary fibration (described by a cohomology class in
$H^{n+2}(X_n,\pi_{n+1})$ ) comes from a``schematic'' one, and that the latter is
essentially unique; in other words, \uline{that the canonical map from the}
\uline{``schematic''} $H^{n+2}$ (\uline{with``quasi-coherent''coefficients}) \uline{to the usual}
\uline{``discrete''} \uline{one is bijective}. Maybe this hope is wholly unrealistic\hand{\unsure{,}}
though. One fact which calls for some scepticism about this approach,
comes in when looking at the case of an ``abelian'' homotopy type, descri-%
% NOTE: the right edge of the scan cuts this line after "descri"
bed by a semisimplicial abelian group $X_{\hand{\ast}}$ , in which case we expect that
% NOTE: the right edge of the scan cuts the last letters of "that"
base change $\underline{Z}\rightarrow k$ should be just the usual base change
\[ X_{\hand{\ast}} \rightarrow X_{\hand{\ast}}\otimes_{\underline{Z}} k \]
(if X. has torsion-free \xout{components} i.e. flat components,at any rate).
But when X. has homology torsion, the universal coefficients formula
shows us that the homology \deleted{(=}\inserted{(or rather,} homotopy) groups of $X_{\hand{\ast}}\otimes_{\underline{Z}}k$ are \uline{not} just
the groups $\unsure{H}_{\unsure{i}}(X_{\hand{\ast}})\otimes_{\underline{Z}}k = \hand{\pi}_i(X_{\hand{\ast}})\otimes_{\underline{Z}}k$ , as we implicitly were assuming it
seems in the approach sketched above, when schematizing the homotopy
groups $\hand{\pi}_i$ \xout{by} one by one via $W(\hand{\pi}_i)$. Thus \inserted{maybe} Postnikov dévissage \inserted{is'nt} \xout{may not be}
a possible approach towards schematization of homotopy types, and one
will have to work out rather a comprehensive yoga of reconstructing
a homotopy type from \xout{its abelianization} one kind or other of``abelianiza-%
tion'' or ``linearization'' \inserted{of homotopy types,} endowed with suitable extra structure embody-%
ing ``multiplicative''\xout{\unsure{f}} features of the homology and cohomology structure.
At any rate, I did not hit upon any ``simplistic way'' to define homotopy
types over\deleted{\illegible{1}} any ground ring, and I have some doubts there is any, in
terms of the general non-sense \xout{of} we did so far\unsure{,|.} \xout{at any rate}

\margin{2}
